% Einheit 1 von 2 (Katalog-Einheit 3): Streifen- und Kreisdiagramm
\einheitenkopf[e1][1 Streifen- und Kreisdiagramm]{Einheit 1 von 2 · Streifen- und Kreisdiagramm}
\zweigzeile{Hier lernst du, Anteile als Streifen und als Kreisdiagramm darzustellen · neu oder schon bekannt – je nach Buch (am Gymnasium schon seit Klasse 6) · P10 oft · baut auf: Prozentsatz (Nr. 1–2), relative Häufigkeit (Nr. 3), Winkel (Nr. 5–6), Längen (Nr. 7)}
\setcounter{aufgabe}{11}

\begin{aufgabe}{Ich erkenne, was das Ganze ist.}
Unterstreiche die Zahl, die das Ganze angibt.
\begin{teile}
\teil Von den 28 Kindern der 7a fahren 12 mit dem Rad.
\teil 45 von 300 Befragten lesen täglich Zeitung.
\teil Ein Verein hat 150 Mitglieder, 60 davon sind Kinder.
\teil 18 Kinder der 7b haben Geschwister. Die Klasse hat 26 Kinder.
\teil Im Kino sahen 90 Gäste den Film A, im ganzen Kino waren 240 Gäste.
\end{teile}
\end{aufgabe}

\begin{aufgabe}{Ich erkenne, ob eine Zahl ein Anteil oder ein Winkel ist.}
Zu einem Kreisdiagramm „Schulweg“ notiert Lisa Zahlen ohne Einheit. Kreuze an.
\begin{teile}
\teil Alle Sektoren zusammen: 360 \hfill \kreuz{Anteil (\%)}\quad\kreuz{Winkel (°)}
\teil Rad: 35 von 100 Kindern \hfill \kreuz{Anteil (\%)}\quad\kreuz{Winkel (°)}
\teil Der Sektor „Bus“ ist ein rechter Winkel: 90 \hfill \kreuz{Anteil (\%)}\quad\kreuz{Winkel (°)}
\teil Alle Anteile zusammen: 100 \hfill \kreuz{Anteil (\%)}\quad\kreuz{Winkel (°)}
\teil Der Sektor „Auto“ ist 54 weit geöffnet. \hfill \kreuz{Anteil (\%)}\quad\kreuz{Winkel (°)}
\end{teile}
\end{aufgabe}

\verfahren{Streifendiagramm}
\begin{aufgabe}{Ich kann einen Anteil als Streifenabschnitt zeichnen.}
Der Streifen ist 10 cm lang, das sind 100\,\%. Färbe den Anteil.
\begin{teile}
\teil $30\,\%$ \quad \streifen{0}{$0\,\%$}{$100\,\%$}
\teil $60\,\%$ \quad \streifen{0}{$0\,\%$}{$100\,\%$}
\teil $90\,\%$ \quad \streifen{0}{$0\,\%$}{$100\,\%$}
\teil $20\,\%$ \quad \streifen{0}{$0\,\%$}{$100\,\%$}
\anweisung{Jetzt ohne Teilstriche: Miss mit dem Lineal.}
\teil $45\,\%$ \quad \streifen[0]{0}{$0\,\%$}{$100\,\%$}
\teil $8\,\%$ \quad \streifen[0]{0}{$0\,\%$}{$100\,\%$}
\end{teile}
\end{aufgabe}

\begin{aufgabe}{Ich kann ein Streifendiagramm mit mehreren Abschnitten zeichnen.}
Zeichne ein Streifendiagramm (10 cm) und beschrifte jeden Abschnitt.
\begin{teile}
\teil „Hast du ein Haustier?“: ja 70\,\%, nein 30\,\%\\ \streifen[0]{0}{$0\,\%$}{$100\,\%$}
\teil Lieblingsobst: Äpfel 45\,\%, Bananen 35\,\%, Birnen 20\,\%\\ \streifen[0]{0}{$0\,\%$}{$100\,\%$}
\teil Schulweg: Rad 35\,\%, Bus 40\,\%, der Rest geht zu Fuß.\\ \streifen[0]{0}{$0\,\%$}{$100\,\%$}
\teil Der Strom einer Gemeinde stammt zu 45\,\% aus Wind, zu 25\,\% aus Sonne, zu 5\,\% aus Wasser und zu 25\,\% aus sonstigen Quellen. Stelle die Anteile in einem 10 cm langen Streifendiagramm dar und beschrifte es. (P10 2018 OS)\\ \streifen[0]{0}{$0\,\%$}{$100\,\%$}
\end{teile}
\end{aufgabe}

\verfahren{Kreisdiagramm}
\begin{aufgabe}{Ich kann den Mittelpunktswinkel zu einem Anteil berechnen.}
Berechne den Mittelpunktswinkel.
\rechenplatz{3}
\begin{teilezwei}
\tz $50\,\%$: \leerfeld[°] & \tz $25\,\%$: \leerfeld[°] \\
\tz $10\,\%$: \leerfeld[°] & \tz $75\,\%$: \leerfeld[°] \\
\tz $15\,\%$: \leerfeld[°] & \tz $5\,\%$: \leerfeld[°] \\
\tz $\frac{3}{8}$ aller Stimmen: \leerfeld[°] & \tz 27 von 120: \leerfeld[°] \\
\tz 14 von 400: \leerfeld[°] & \\
\anweisung{Runde auf eine Stelle nach dem Komma.}
\tz 5 von 13: \leerfeld[°] & \\
\end{teilezwei}
\end{aufgabe}

\begin{aufgabe}{Ich kann einen Sektor ab dem Radius zeichnen.}
Trage den Winkel am Radius ab und färbe den Sektor.

\noindent\begin{minipage}[t]{0.32\linewidth}
\begin{teile}\teil 90°\end{teile}
\kreisleer[1.8]
\end{minipage}\hfill
\begin{minipage}[t]{0.32\linewidth}
\begin{teile}\teil 45°\end{teile}
\kreisleer[1.8]
\end{minipage}\hfill
\begin{minipage}[t]{0.32\linewidth}
\begin{teile}\teil 40\,\% – berechne zuerst den Winkel.\end{teile}
\kreisleer[1.8]
\end{minipage}
\end{aufgabe}

\begin{aufgabe}{Ich kann ein Kreisdiagramm zeichnen. (kennst du wahrscheinlich seit Klasse 6)}
Zeichne das Kreisdiagramm und beschrifte die Sektoren.

\noindent\begin{minipage}[t]{0.32\linewidth}
\begin{teile}\teil Rad 50\,\%, Bus 25\,\%, zu Fuß 25\,\%\end{teile}
\kreisleer[1.8]
\end{minipage}\hfill
\begin{minipage}[t]{0.32\linewidth}
\begin{teile}\teil Sport 55\,\%, Musik 30\,\%, Lesen 15\,\%\end{teile}
\kreisleer[1.8]
\end{minipage}\hfill
\begin{minipage}[t]{0.32\linewidth}
\begin{teile}\teil Von 24 Kindern spielen 12 Fußball, 8 Handball und 4 Basketball.\end{teile}
\kreisleer[1.8]
\end{minipage}
\end{aufgabe}

\begin{aufgabe}{Ich kann den Anteil eines Sektors schätzen.}
Welcher Anteil vom ganzen Kreis ist grau? Gib ihn als Bruch an.

\noindent\begin{minipage}[t]{0.32\linewidth}
\begin{teile}\teil \leerfeld\end{teile}
\kreissektor{180}{}
\end{minipage}\hfill
\begin{minipage}[t]{0.32\linewidth}
\begin{teile}\teil \leerfeld\end{teile}
\kreissektor{90}{}
\end{minipage}\hfill
\begin{minipage}[t]{0.32\linewidth}
\begin{teile}\teil \leerfeld\end{teile}
\kreissektor{120}{}
\end{minipage}

\noindent\begin{minipage}[t]{0.32\linewidth}
\begin{teile}\teil \leerfeld\end{teile}
\kreissektor{60}{}
\end{minipage}\hfill
\begin{minipage}[t]{0.64\linewidth}
\anweisung{Gib den Anteil in Prozent an.}
\begin{teile}\teil \leerfeld[\%]\end{teile}
\kreissektor{144}{}
\end{minipage}
\end{aufgabe}

\begin{aufgabe}{Ich kann Sektoren nach ihrer Größe zuordnen.}
Schreibe die Angaben an die passenden Sektoren.

\noindent\begin{minipage}[t]{0.48\linewidth}
\begin{teile}\teil Rad 60\,\%, Bus 25\,\%, zu Fuß 15\,\%\end{teile}
\kreisdiagrammleer{Bus/25, zu Fuß/15, Rad/60}
\end{minipage}\hfill
\begin{minipage}[t]{0.48\linewidth}
\begin{teile}\teil Lieblingsfarbe: Blau 35\,\%, Rot 30\,\%, Grün 20\,\%, Gelb 15\,\%\end{teile}
\kreisdiagrammleer{Grün/20, Blau/35, Gelb/15, Rot/30}
\end{minipage}
\end{aufgabe}

\begin{aufgabe}{Ich kann Sektoren nach ihrer Größe zuordnen – weiter.}
Schreibe die Angaben an die passenden Sektoren.

\noindent\begin{minipage}[t]{0.40\linewidth}
\begin{teile}\teil Getränke beim Schulfest: Tee 26\,\%, Kaffee 24\,\%, Saft 30\,\%, Wasser 20\,\%\end{teile}
\kreisdiagrammleer{Kaffee/24, Saft/30, Wasser/20, Tee/26}
\end{minipage}\hfill
\begin{minipage}[t]{0.56\linewidth}
\begin{teile}\teil Ein Haushalt verbraucht an einem Tag 490 Liter Wasser: Duschen und Baden 175 l, Toilette 130 l, Wäsche waschen 72 l, Kochen und Trinken 17 l, Sonstiges 96 l. Ordne die fünf Bereiche den Sektoren zu und berechne den Mittelpunktswinkel für Kochen und Trinken. Runde auf eine Stelle. (P10 2024 OS)\end{teile}
\kreisdiagrammleer{Sonstiges/96, Duschen/175, Kochen/17, Toilette/130, Wäsche/72}
\end{minipage}
\end{aufgabe}

\begin{aufgabe}{Ich finde den Fehler beim Kreisdiagramm.}
Finde den Fehler, benenne ihn und korrigiere die Rechnung.
\begin{teile}
\teil Bei einer Umfrage nennen 20\,\% den Bus. Tim berechnet den Mittelpunktswinkel für „Bus“:
\rechnung{\text{Winkel Bus} &= 20°}
\par\vspace{16mm}
\teil Von 400 Befragten hören 15\,\% am liebsten Radio. Ben berechnet den Mittelpunktswinkel für „Radio“:
\rechnung{15\,\% \text{ von } 400 &= 60 \\ \text{Winkel Radio} &= 60°}
\par\vspace{16mm}
\end{teile}
\end{aufgabe}

\begin{aufgabe}{Ich kann begründen, warum 1\,\% im Kreisdiagramm 3,6° sind.}
Begründe ohne zu zeichnen.
\begin{teile}
\teil Warum entspricht 1\,\% genau einem Winkel von 3,6°?
\par\vspace{12mm}
\teil Kann ein Kreisdiagramm die Anteile 45\,\%, 35\,\% und 30\,\% zeigen?
\par\vspace{12mm}
\end{teile}
\end{aufgabe}

\begin{aufgabe}{Ich kann Umfrageergebnisse in einem Kreisdiagramm darstellen.}
Die Schülervertretung (SV) hat 240 Schülerinnen und Schüler gefragt, wohin der Wandertag gehen soll.

\sachtabelle{lccc}{Ziel & Stimmen & Anteil & Mittelpunktswinkel}{Kletterpark & 96 & \leerzelle & \leerzelle\\ Zoo & 60 & \leerzelle & \leerzelle\\ Schwimmbad & 51 & \leerzelle & \leerzelle\\ Museum & 33 & \leerzelle & \leerzelle}

\noindent\begin{minipage}[t]{0.55\linewidth}
\begin{teile}
\teil Berechne die Anteile in Prozent und die Mittelpunktswinkel. Trage sie in die Tabelle ein.
\teil Zeichne das Kreisdiagramm.
\teil Die SV bietet nur Ziele an, die mindestens ein Viertel der Stimmen haben. Welche Ziele bleiben? Begründe.
\end{teile}
\end{minipage}\hfill
\begin{minipage}[t]{0.42\linewidth}
\kreisleer[2.2]
\end{minipage}
\end{aufgabe}
